% Generated draft: requires mathematical and source review.
\documentclass[11pt]{article}
\usepackage[margin=25mm]{geometry}
\usepackage{fontspec}
\setmainfont{Latin Modern Roman}
\usepackage{amsmath,amssymb,mathtools,unicode-math}
\setmathfont{Latin Modern Math}
\usepackage{enumitem,needspace,xurl,hyperref}
\setlength{\parindent}{0pt}
\setlength{\parskip}{5pt}
\setlength{\emergencystretch}{3em}
\newcommand{\N}{\mathbb{N}}
\newcommand{\Z}{\mathbb{Z}}
\newcommand{\Q}{\mathbb{Q}}
\newcommand{\R}{\mathbb{R}}
\newcommand{\C}{\mathbb{C}}
\newcommand{\F}{\mathbb{F}}
\newcommand{\A}{\mathbb{A}}
\newcommand{\PP}{\mathbb{P}}
\newcommand{\E}{\mathbb{E}}
\newcommand{\eps}{\varepsilon}
\newcommand{\dd}{\,\mathrm{d}}
\begin{document}
\section*{UnsolvedMath ID: 701601}
\textbf{Source problem code:} catalogue entry
\subsection{Definitions and mathematical statement}
For a group \(G\), let \(1_G\) denote its identity and let \(\langle x,y,z\rangle\) denote the subgroup generated by \(x,y,z\). The condition \(g^{12}=1_G\) for every \(g\in G\) means that every element has finite order dividing \(12\); it does not require every element to have order exactly \(12\).

Is the following assertion true? For every group \(G\) and every \(x,y,z\in G\), if
\[
G=\langle x,y,z\rangle,\qquad
x^3=y^2=z^2=(xy)^3=(yz)^3=1_G,
\]
and
\[
\forall g\in G,\qquad g^{12}=1_G,
\]
then
\[
|G|\leq 12.
\]
No finiteness assumption is imposed on \(G\). The displayed relations require the relevant orders to divide the displayed exponents, rather than to equal them.

The supplied source states that the conclusion holds when \(G\) is finite, attributing this result to D. V. Lytkina and V. D. Mazurov (2015). This finite-case result is retained as a supplied literature qualification, not as an independently verified assertion.
\subsection{Short English statement}
Must every group generated by three elements satisfying the displayed relations, in which every element has order dividing \(12\), have at most \(12\) elements? The supplied record labels the problem open and reports that the bound is known for finite groups.
\subsection{Sources}
\begin{itemize}
\item Local UnsolvedMath registry, supplied record with permanent ID \(701601\), exact source problem code \texttt{catalogue entry}, titled \emph{Kourovka Notebook Problem 19.53}. The supplied status is \texttt{open}; this label has not been independently verified.
\item \emph{The Kourovka Notebook (Unsolved Problems in Group Theory)}, Problem 19.53, proposed by V. D. Mazurov in 2018. The supplied background identifies the 19th issue (2018) and the 21st issue edition (2026), arXiv:1401.0300v48. The supplied collection description names E. I. Khukhro and V. D. Mazurov as editors.
\item D. V. Lytkina and V. D. Mazurov, \emph{Siberian Math. J.}, \textbf{56}, no. 3 (2015), 471--475. The supplied statement cites this reference for the affirmative finite-group case; no article title or URL is supplied.
\end{itemize}
\end{document}
